\chapter{Testes e exercicios}

\section{Por que testar com XML pequeno?}

Voce nao deve testar com o arquivo gigante. Testes devem ser pequenos, rapidos e previsiveis.

Crie:

\begin{verbatim}
tests/fixtures/small_export.xml
\end{verbatim}

Conteudo exemplo:

\begin{lstlisting}[style=devbutterbash]
<HealthData>
    <ExportDate value="2026-06-23 08:00:00 -0300" />
    <Record
        type="HKQuantityTypeIdentifierHeartRate"
        sourceName="Apple Watch"
        unit="count/min"
        startDate="2026-06-23 07:00:00 -0300"
        endDate="2026-06-23 07:00:00 -0300"
        value="90"
    />
    <Record
        type="HKCategoryTypeIdentifierSleepAnalysis"
        sourceName="Apple Watch"
        startDate="2026-06-23 00:00:00 -0300"
        endDate="2026-06-23 07:00:00 -0300"
        value="HKCategoryValueSleepAnalysisAsleepCore"
    />
</HealthData>
\end{lstlisting}

\section{Testes de funcoes puras}

\begin{lstlisting}[style=devbutterpython]
from src.parser.xml_parser import clean_tag, safe_float


def test_clean_tag_without_namespace():
    assert clean_tag("Record") == "Record"


def test_clean_tag_with_namespace():
    tag = "{http://www.topografix.com/GPX/1/1}trkpt"
    assert clean_tag(tag) == "trkpt"


def test_safe_float_number():
    assert safe_float("90") == 90.0


def test_safe_float_text():
    assert safe_float("HKCategoryValueSleepAnalysisAsleepCore") is None
\end{lstlisting}

\section{Exercicios progressivos}

\subsection{Exercicio 1: Counter}

Conte palavras:

\begin{lstlisting}[style=devbutterpython]
words = ["heart", "steps", "heart", "sleep"]
\end{lstlisting}

Resultado esperado:

\begin{lstlisting}[style=devbutterpython]
{"heart": 2, "steps": 1, "sleep": 1}
\end{lstlisting}

Depois aplique a mesma ideia para contar tipos de \code{Record}.

\subsection{Exercicio 2: batch}

Crie uma funcao que recebe numeros de 1 a 20 e separa em batches de 5.

Entrada:

\begin{lstlisting}[style=devbutterpython]
list(range(1, 21))
\end{lstlisting}

Saida:

\begin{lstlisting}[style=devbutterpython]
[[1, 2, 3, 4, 5], [6, 7, 8, 9, 10], ...]
\end{lstlisting}

\subsection{Exercicio 3: queue}

Crie um producer que coloca numeros em uma fila e um consumer que imprime os numeros.

\subsection{Exercicio 4: threading}

Rode producer e consumer em threads separadas e coloque o nome da thread no log.

\subsection{Exercicio 5: profiling}

Meça o tempo de uma funcao usando \code{perf_counter}.

\subsection{Exercicio 6: banco}

Insira 1000 registros fake no Postgres com:

\begin{itemize}
  \item insert linha por linha;
  \item batch insert;
\end{itemize}

Compare o tempo.

\section{Testes de banco}

No comeco, pode testar manualmente. Depois, use um banco de teste via Docker.

Teste importante:

\begin{itemize}
  \item inserir batch;
  \item contar linhas;
  \item consultar por tipo;
  \item garantir que \code{value_numeric} vira \code{NULL} quando valor e texto.
\end{itemize}
